\documentclass[12pt, letterpaper]{article}

%-------Packages---------
\usepackage{amssymb,amsfonts}
\usepackage{mathptmx}
\usepackage{amsthm}
\usepackage{graphicx}
\usepackage[all,arc]{xy}
\usepackage[colorlinks=true, allcolors=blue]{hyperref}
\usepackage{enumerate}
\usepackage{bbm}
\usepackage{dsfont}
\usepackage{mathrsfs}
\usepackage{tikz-cd}
\usepackage{setspace}
\usepackage{import}
\usepackage[utf8]{inputenc}
\usepackage{hyperref}
\usepackage{tikz}
\usepackage{float}
\usepackage{mdframed}
\usepackage{fancyhdr}
\usepackage[nointegrals]{wasysym}

\usepackage[left=1in,top=1in,right=1in,bottom=1in]{geometry}

\usepackage{mathtools}
\usepackage{setspace}
\usepackage{cleveref}
\usepackage{multicol}

%--------Theorem Environments--------
%theoremstyle{plain} --- default
\newmdtheoremenv{theorem}{Theorem}[subsection]
\newmdtheoremenv{corollary}[theorem]{Corollary}
\newtheorem{proposition}[theorem]{Proposition}
\newmdtheoremenv{lemma}[theorem]{Lemma}
\newtheorem{conj}[theorem]{Conjecture}
\newtheorem{quest}[theorem]{Question}

\theoremstyle{definition}
\newmdtheoremenv{definition}[theorem]{Definition}
\newmdtheoremenv{definitions}[theorem]{Definitions}
\newtheorem{example}[theorem]{Example}
\newtheorem{algorithm}[theorem]{Algorithm}
\newtheorem{rmk}[theorem]{Remark}
\newtheorem{exercise}[theorem]{Exercise}

%----------- my commands -------------
\newcommand{\C}{\mathbb{C}}
\newcommand{\R}{\mathbb{R}}
\newcommand{\Q}{\mathbb{Q}}
\newcommand{\Z}{\mathbb{Z}}
\newcommand{\N}{\mathbb{N}}
\renewcommand{\P}{\mathbb{P}}
\newcommand{\contr}{\Rightarrow \Leftarrow}
\newcommand{\HRule}[1]{\rule{\linewidth}{#1}}
\newcommand{\s}{\(\,\)}
\renewcommand{\t}[1]{\text{#1}}
\newcommand{\ang}[1]{\left\langle #1 \right\rangle}

% Meta data
\setstretch{1.2}

\begin{document}

\pagestyle{fancy}
\fancyhf{}
\fancyhead[LOE]{\slshape{\textbf{Vincent Ferrara}}}
\fancyhead[ROE]{\thepage}

\section*{HW 1}

\subsection*{Problem 1}
Let $G$ be a group and $N \unlhd G$\\
Let $\psi$ be an homomorphism form $G$ to $G_1$ s.t. $N \subseteq \text{ ker}(\psi)$\\
Define: $\bar{\psi}: G \diagup N \rightarrow G_1$ s.t. $\bar{\psi}(gN)=\psi(g)$\\\
Well Defined: Suppose $gN=hN \implies gh^{-1} \in N \implies gh^{-1} \in $ ker$(\psi)$\\
$\psi(gh^{-1})=e \implies \psi(g)=\psi(h)$ since $\psi$ is a homomorphism this implies $\bar{\psi}(gN)=\bar{\psi}(hN)$\\
Homomorphism: Let $gN,hN \in G \diagup N$\\
$\bar{\psi}(gNhN)=\bar{\psi}(ghN)=\psi(gh)=\psi(g)\psi(h)=\bar{\psi}(gN)\bar{\psi}(hN)$\\
Uniqueness: Suppose there exists another homomorphism $\phi: G\diagup  N \rightarrow G_1$ such that $\varphi=\psi \circ \varphi$\\
\[
\phi(gN)=\psi(g) \implies \phi = \bar{\psi}
\]

\newpage
\subsection*{Problem 2}
Let $H$ be a subgroup of $G$.
\begin{itemize}
    \item (Inverses) Let $a \in N_G(H)$\\
    Consider: $a^{-1}Ha=a^{-1}aH=H$ thus $a^{-1} \in N_G(H)$
    \item (Closure) Let $a,b \in N_G(H)$\\
    Consider: $abHb^{-1}a^{-1} = abb^{-1}Ha^{-1}=aHa^{-1}=H$ thus $ab \in N_G(H)$
    \item (Identity) Consider $eHe=H$ thus $e \in N_G(H)$
\end{itemize}
Let $h \in H$\\
Consider : $hHh^{-1}$ since $h,h^{-1} \in H \implies hHh^{-1}=H$ thus $h \in N_G(H)$ thus $H \subseteq N_G(H)$\\
By definition since $\forall g \in N_g(H)\; gHg^{-1}=H$ thus $H \unlhd N_G(H)$\\
Largest: Let $K$ be a subgroup of $G$ s.t. $H \unlhd K$ thus $\forall k \in K$ 
\[kHk^{-1}=H \]
thus $k \in N_G(H)$ thus $K \subseteq N_G(H)$

\newpage
\subsection*{Problem 3}
\begin{itemize}
    \item (i. $\implies$ ii)
    \begin{proof}
        Assume i.\\
        Let $g \in H \cap K \implies g \in H$ and $g \in K$
    \end{proof}
\end{itemize}
\end{document}